% Toward Gauge Structure from Non-Injective Projection
% Cosmochrony companion paper M4a

\documentclass[11pt,a4paper]{article}

\usepackage[T1]{fontenc}
\usepackage{lmodern}
\input{glyphtounicode}
\pdfgentounicode=1
\usepackage{amsmath,amssymb,amsthm}
\usepackage{enumerate}
\usepackage{hyperref}
\hypersetup{pdftitle={Toward Gauge Structure from Non-Injective Projection: Fibre-Invariant Charges, a Supplied SU(2)
  Carrier, and a Conditional Connection}, pdfauthor={J\'er\^ome Beau},
  pdfsubject={Fibre-invariant charges, a supplied SU(2) carrier and a conditional gauge connection from non-injective
  projection},
  pdfkeywords={gauge symmetry; non-injective projection; fibre bundles; gauge invariance; Yang\textendash Mills theory;
  emergent gauge structure; admissibility; Heisenberg group; SU(2); quaternionic structure}}
\usepackage{doi}

\newtheorem{theorem}{Theorem}[section]
\newtheorem{lemma}[theorem]{Lemma}
\newtheorem{proposition}[theorem]{Proposition}
\newtheorem{corollary}[theorem]{Corollary}
\newtheorem{definition}[theorem]{Definition}
\newtheorem{remark}[theorem]{Remark}

\title{Toward Gauge Structure from Non-Injective Projection:\\
Fibre-Invariant Charges, a Supplied SU(2) Carrier, and a Conditional Connection}

\author{J\'er\^ome Beau\\
\small Independent Researcher\\
\small \texttt{jerome.beau@cosmochrony.org}}

\date{2026}

\begin{document}

  \maketitle

  \begin{abstract}
    Starting from the axioms of admissible non-injective transitions of the Foundation paper, we
    state which gauge-theoretic structures carried by the fibres of a non-injective projection are
    defined, which are conditional, and which remain open.
    Gauge charges are functions constant on the projection fibres; they factor through the
    projection, so they are functions of the observable state, and their identification with
    conserved physical charges is open.
    The gauge group is the group of fibre-preserving transformations.
    $\mathrm{U}(1)$ is the residual global phase symmetry of span-based observables.
    On the spinor carrier $V_\rho \cong \mathbb{C}^2$ supplied conditionally by paper O23, and
    under the neutral-sector hypothesis of paper O27, the group acting faithfully on the carrier is
    $\mathrm{SU}(2)$; this is a consistency statement about a supplied carrier, not a selection.
    $\mathrm{SU}(3)$ is a structural hypothesis.
    Two unproved continuum hypotheses, a smooth fibre bundle and a differentiable base, the second
    the open spatial continuum question of paper Q5a, do not determine a connection.
    On a supplied structure group, representation and connection, the curvature measures the path
    dependence of transport around infinitesimal loops.
    The Yang--Mills action is a structural motivation on this curvature, not a derivation.
    Interpretive outlook, a reading and not a result: gauge redundancy is read as the symmetry of a
    fibre that the observables do not resolve, and the paper locates where the group, the
    comparison of fibres and the dynamics would enter, without supplying them.
  \end{abstract}

  \paragraph{Keywords.}
    gauge symmetry; non-injective projection; fibre bundles; gauge invariance; Yang–Mills theory; emergent gauge
    structure; admissibility; Heisenberg group; SU(2); quaternionic structure

    \clearpage
    \tableofcontents

    \clearpage
  \section{Introduction}
  \label{sec:intro}

  \subsection*{Position in the programme}

    The Foundation paper~\cite{Beau2026m} states the axioms A1--A4 of admissible non-injective
    transitions and derives from A1 and A3 the structural indeterminacy of the proto-state, an unresolved
    configuration open across the admissible directions of the fibre.
    A3 preserves those open directions but does not by itself forbid an erasure of their relative phases,
    so the retention of phase coherence is not derived from A3 alone.
    These results concern the \emph{internal} structure of the admissible fibre at a
    single transition step.

    The present paper addresses one macroscopic consequence of allowing the projection
    structure to vary from point to point in the effective continuum: the invariants of the
    projection fibre and the group of transformations that preserves the fibres.
    The answer is gauge charges and a gauge group (§\ref{sec:charges}), with a connection
    and its curvature defined on supplied transport data under stated continuum hypotheses
    (§\ref{sec:connection}--\ref{sec:ym}).
    Effective spacetime geometry is not addressed: the spatial continuum it requires is an open
    question of~\cite{Beau2026q5a}.

  \subsection*{What is and is not derived in this paper}

    The paper defines gauge charges and the gauge group from the fibre structure of a
    non-injective projection, and identifies the residual phase symmetry $\mathrm{U}(1)$ when the projection
    is defined through span-based observables.
    It does not select the gauge group.
    Its non-abelian part is read off a carrier that is supplied, not derived: on the spinor
    carrier of~\cite{Beau2026a27}, and under the neutral-sector hypothesis of~\cite{Beau2026a31},
    it is $\mathrm{SU}(2)$ (Proposition~\ref{thm:su2-rigidity}).
    $\mathrm{SU}(3)$ is a hypothesis.
    Gauge charges are functions of the observable state; their identification with conserved
    physical charges is not established.
    The connection is not derived: the two continuum hypotheses that the paper states but does not
    prove are needed for a continuum description and do not determine a connection, so a structure
    group, its representation and a connection are supplied (§\ref{subsec:transport}).
    The curvature and the Yang\textendash Mills action are statements on these supplied data, and
    the Yang\textendash Mills action is a structural motivation, not a derivation.
    Coupling constants and the matter content of the Standard Model are not derived.
    These points are collected in Section~\ref{sec:open}.

  \subsection*{Organisation}

    Section~\ref{sec:charges} defines gauge charges as fibre invariants and the gauge group as
    the group of fibre-preserving transformations, and states what fixes its non-abelian part
    on a supplied carrier.
    Section~\ref{sec:connection} shows that the continuum hypotheses do not fix a connection, states
    the supplied transport data, and defines on them the gauge connection as the infinitesimal
    generator of parallel transport.
    Section~\ref{sec:ym} gives the curvature of this connection, its local and global meaning, and
    the structural motivation for the Yang\textendash Mills action.
    Section~\ref{sec:open} identifies the open directions, and Section~\ref{sec:conclusion}
    concludes with a labelled interpretive outlook.

    \clearpage
  \section{Gauge Charges as Fibre Invariants}
  \label{sec:charges}

  \subsection{Charges from non-injectivity}

    The non-injectivity of $\Pi_n$ (Axiom~A2 of~\cite{Beau2026m}) means that
    multiple configurations in $\Omega_n^{(c)}$ project to the same observable
    state $O_n$.
    These configurations are physically indistinguishable at the observable level
    but may carry distinct internal structure.

    \begin{definition}[Gauge charge]
      \label{def:gauge-charge}
      A \emph{gauge charge} is a functional
      \begin{equation}
        Q : \Omega_n^{(c)} \;\longrightarrow\; \mathbb{R}
      \end{equation}
      that is constant on projection fibres:
      \begin{equation}
        \Pi_n(\omega_1) = \Pi_n(\omega_2)
        \;\Longrightarrow\;
        Q(\omega_1) = Q(\omega_2).
      \end{equation}
    \end{definition}

    Every gauge charge factors through the projection.
    On the image of $\Pi_n$, set $q(O) := Q(\omega)$ for any $\omega$ with $\Pi_n(\omega) = O$; the
    definition makes $q$ independent of the chosen preimage, so $Q = q \circ \Pi_n$.
    A gauge charge is therefore a function of the observable state: it is resolved by $\Pi_n$, and it
    carries no distinction within a fibre.
    What non-injectivity adds is the converse restriction: a functional that separates two points of
    one fibre is not a gauge charge.
    Fibre constancy is an invariance property.
    It does not by itself supply a conservation law, a charge generator, a matter representation, or
    an identification with electric, weak or colour charge; these remain open (Section~\ref{sec:open}).

    \begin{definition}[Gauge group]
      \label{def:gauge-group}
      The \emph{gauge group} $\mathcal{G}$ is the group of transformations
      $\omega \mapsto g \cdot \omega$ that preserve the projection:
      \begin{equation}
        \mathcal{G}
        \;:=\;
        \bigl\{g \in \mathrm{Aut}(\Omega_n^{(c)}) \mid
        \Pi_n(g \cdot \omega) = \Pi_n(\omega) \text{ for all } \omega \bigr\}.
      \end{equation}
    \end{definition}

    $\mathcal{G}$ consists precisely of the fibre-preserving symmetries: it acts
    within $\Pi_n^{-1}(O_n)$ without affecting the observable output.
    Its orbits lie inside the fibres, and they exhaust a fibre only if $\mathcal{G}$ acts transitively
    on it, for instance when $\mathrm{Aut}(\Omega_n^{(c)})$ contains every permutation of the points of
    that fibre.
    Equality of observable output and equivalence under $\mathcal{G}$ therefore coincide only under
    such a transitivity condition.

    \begin{remark}[Gauge invariance as structural necessity]
        The gauge group is not introduced by hand.
        It is fixed by the structure of $\Pi_n$ once $\mathrm{Aut}(\Omega_n^{(c)})$ is specified.
        Provided $\mathrm{Aut}(\Omega_n^{(c)})$ contains the permutations of points within a
        fibre, a non-injective $\Pi_n$ has non-trivial fibre-preserving transformations, and
        hence a non-trivial gauge group.
        Under that proviso, gauge invariance is a structural consequence of Axiom~A2, not an
        additional postulate.
    \end{remark}

  \subsection{Structural constraints on admissible gauge groups}
    \label{subsec:gauge-constraints}

    Definition~\ref{def:gauge-group} alone does not restrict $\mathcal{G}$ to a particular
    Lie group.
    Three structural inputs bear on it; each is stated with its status.

    \begin{enumerate}[(i)]
        \item \emph{Non-factorisability.}
        Once a representation-theoretic dictionary for $F_n$ is supplied, Axioms~A1 and~A3 make the
        admissible generators act irreducibly on $F_n$ only under an additional sector reading of A3,
        not derived, on which a decomposition of $F_n$ stable under those generators counts as premature
        selection (Foundation Lemma~5.6 of~\cite{Beau2026m}).
        This constrains the admissible generators, not the fibre-preserving group
        $\mathcal{G}$ directly; under that reading, it is the structural reason why the fibre is not treated
        as a product of independent components.

        \item \emph{A supplied carrier.}
        The axioms do not identify $F_n$ with a Heisenberg or Weil carrier (Foundation
        Theorem~5.7 of~\cite{Beau2026m}).
        The spectral programme supplies a Heisenberg carrier, on which the fibre carries the
        Heisenberg\textendash Schr\"odinger representation of
        $\mathrm{Heis}_3(\mathbb{Z}/q\mathbb{Z})$ at a fixed central character~\cite{Beau2026q5a}.
        Within that supplied setting, O23~\cite{Beau2026a27} (Theorem~3.1) shows that if the
        admissible neutral sector is carried by an irreducible two-dimensional
        $\mathrm{SU}(2)$-valued $V_\rho \cong \mathbb{C}^2$, its traceless neutral sector is
        $\mathfrak{su}(2) \cong \mathrm{Im}\,\mathbb{H}$, of real dimension three.
        The selection of that carrier is open in O23.

        \item \emph{Stability under admissibility refinement.}
        Under spectral refinement (large-$q$ limit), admissible sectors are required to remain
        closed under the projection-induced transport.
        This working requirement motivates restricting attention to compact Lie groups; it is
        not derived here.
    \end{enumerate}

    These inputs do not determine the gauge group.
    $\mathrm{U}(1)$ is the residual phase symmetry when $\Pi_n$ is span-based (§\ref{subsec:identification}); the
    non-abelian part is $\mathrm{SU}(2)$ only on the supplied carrier of input~(ii)
    (Proposition~\ref{thm:su2-rigidity}); $\mathrm{SU}(3)$ remains a hypothesis.

  \subsection{Identification of the gauge groups}
    \label{subsec:identification}

    The three groups enter with different statuses.

    \begin{itemize}
      \item \emph{$\mathrm{U}(1)$}: global phase rotations
      $\omega \mapsto e^{i\theta}\omega$ within the complex fibre $\mathbb{C}^q$.
      They act as a global scalar, which leaves every span of fibre vectors unchanged, and
      hence every capacity profile defined from such spans~\cite{Beau2026a16}; the global
      phase is therefore not resolved by these observables, while the rotation acts
      non-trivially on the configuration.
      The rotations belong to $\mathcal{G}$ to the extent that $\Pi_n$ is defined through such
      span-based observables.

      \item \emph{$\mathrm{SU}(2)$}: the group acting faithfully on the spinor carrier
      $V_\rho \cong \mathbb{C}^2$ supplied conditionally by O23~\cite{Beau2026a27}, under the
      neutral-sector hypothesis of O27~\cite{Beau2026a31}
      (Proposition~\ref{thm:su2-rigidity}).
      The carrier is an input; its selection is open.

      \item \emph{$\mathrm{SU}(3)$}: a structural hypothesis for a higher-dimensional
      internal symmetry of the fibre, corresponding to several co-admissible internal
      directions beyond the $\mathrm{SU}(2)$ carrier.
      Its identification is an open direction (Section~\ref{sec:open}).
    \end{itemize}

    \begin{remark}[Status of $\mathrm{SU}(3)$]
        The identification of $\mathrm{SU}(3)$ as a gauge group is a structural
        hypothesis, not a derived result.
        Establishing it would require a carrier, supplied or derived, on which a group with
        $\mathrm{SU}(3)$ structure acts within the admissible fibre, together with an
        independent selection argument.
        A spectral degeneracy between representation sectors would not suffice: the exact
        degeneracy $\lambda_{1/2} = \lambda_{3/2}$ on the binary icosahedral group is a fact
        about the group and its generating set and carries no selection content on its
        own~\cite{Beau2026q2}.
    \end{remark}

  \subsection{The non-abelian sector on a supplied carrier}
    \label{subsec:su2-rigidity}

    \begin{proposition}[Non-abelian sector on a supplied carrier]
      \label{thm:su2-rigidity}
      Assume Axioms~A1--A3, together with:
      \begin{enumerate}[(H1)]
        \item the admissible neutral sector is carried by an irreducible two-dimensional
        $\mathrm{SU}(2)$-valued representation $V_\rho \cong \mathbb{C}^2$, as supplied
        conditionally by O23~\cite{Beau2026a27};
        \item the image of the admissible projection is the neutral traceless part of the
        ambient observable space, $\mathrm{Im}\,\Pi_{\mathrm{adm}} = \mathrm{Ntrl}$, as assumed
        in O27~\cite{Beau2026a31}.
      \end{enumerate}
      Then the effective neutral sector $H_{\mathrm{eff}} = \mathrm{Im}\,\Pi_{\mathrm{adm}} \cap
      \mathrm{Ntrl}$ is isomorphic to $\mathfrak{su}(2) \cong \mathrm{Im}\,\mathbb{H}$, and the
      connected compact Lie group with this Lie algebra that acts faithfully on $V_\rho$ is
      $\mathrm{SU}(2)$.
      Hypotheses (H1) and (H2) are load-bearing; neither is a consequence of A1--A3.
    \end{proposition}

    \begin{proof}
        By O23~\cite{Beau2026a27} (Theorem~3.1), under (H1) the traceless part of the neutral
        sector carried by $V_\rho$, inside $\mathrm{End}(V_\rho)$, is
        $\mathfrak{su}(2) \cong \mathrm{Im}\,\mathbb{H}$, of real dimension three.
        By O27~\cite{Beau2026a31}, an admissible morphism satisfying the three constraints of that
        paper factors uniquely and real-linearly through the admissible quotient; under (H2)
        that quotient reaches $H_{\mathrm{eff}}$, and O23's dimension count transfers to it, so
        $H_{\mathrm{eff}} \cong \mathfrak{su}(2)$.
        The connected compact Lie groups with Lie algebra $\mathfrak{su}(2)$ are
        $\mathrm{SU}(2)$ and $\mathrm{SO}(3)$.
        Only $\mathrm{SU}(2)$ acts faithfully on the two-dimensional $V_\rho$, since the
        irreducible two-dimensional representation of $\mathfrak{su}(2)$ does not descend to
        $\mathrm{SO}(3)$.
    \end{proof}

    \begin{remark}[Status of Proposition~\ref{thm:su2-rigidity}]
      Proposition~\ref{thm:su2-rigidity} is a consistency statement: the group is read off a
      carrier that is $\mathrm{SU}(2)$-valued by hypothesis~(H1).
      It does not select $\mathrm{SU}(2)$ among admissible gauge groups, and no result used here
      establishes that $\mathrm{SU}(2)$ is the unique or minimal non-abelian group compatible
      with admissible non-injective projection
      (Remark~\ref{rem:no-quaternionic-forcing}).
      The carrier selection is open in O23~\cite{Beau2026a27}, and hypothesis~(H2) is assumed
      and derived nowhere~\cite{Beau2026a31}.
    \end{remark}

    \begin{remark}[The definition does not force a quaternionic group]
      \label{rem:no-quaternionic-forcing}
      No input above forces the non-abelian part of $\mathcal{G}$ into the unit quaternions.
      If $\mathrm{Aut}(\Omega_n^{(c)})$ contains the permutations of a fibre with at least three
      points, $\mathcal{G}$ contains the symmetric group $\mathfrak{S}_3$, which does not embed
      in $\mathrm{SU}(2)$: the only element of order two in $\mathrm{SU}(2)$ is $-I$, whereas
      $\mathfrak{S}_3$ has three.
      The restriction to $\mathrm{SU}(2)$ in Proposition~\ref{thm:su2-rigidity} is therefore a
      consequence of the supplied carrier, not of Definition~\ref{def:gauge-group}.
    \end{remark}

  \section{The Gauge Connection}
  \label{sec:connection}

  \subsection{The problem of comparing neighbouring fibres}

    The following development is conditional on two structural hypotheses.
    First, the admissible fibre $F_n$ admits an effective continuum description
    in which it forms a smooth fibre bundle over an effective spacetime manifold.
    Second, this manifold itself carries a differentiable structure, which requires
    the emergence of a spatial continuum from the discrete admissibility filter.
    That emergence is open: on a supplied Heisenberg carrier the canonical filtration of the
    fibre is a growing Fourier window, and, under that paper's conditions on the saturation depth and
    the weights, no common scalar normalisation of the published admissibility form yields a
    non-trivial toric differential operator~\cite{Beau2026q5a}.
    Neither hypothesis is an axiom of the present framework or a result of the programme;
    both are working assumptions.
    All results in §\ref{sec:connection}--\ref{sec:ym} are therefore conditional
    on these two hypotheses.
    The two hypotheses are needed for the continuum description below, but they are not sufficient
    for it: they do not determine a connection (§\ref{subsec:transport}).

    For each point $x$ in the effective manifold, the fibre $F_x = \Pi^{-1}(x)$
    carries the internal gauge structure of Section~\ref{sec:charges}.

    When the projection $\Pi(x)$ varies from point to point, the fibres $F_x$ and
    $F_{x+dx}$ are not canonically identified.
    A state $\psi(x)$ transforming under a local gauge transformation $g(x)$, with values in the
    structure group $G$ of §\ref{subsec:transport}, acquires an additional term
    under differentiation:
    \begin{equation}
      \partial_\mu(g(x)\,\psi(x))
      \;=\; g(x)\,\partial_\mu\psi(x) + (\partial_\mu g(x))\,\psi(x),
    \end{equation}
    which is not covariant under local gauge transformations $g(x) \in G$.

  \subsection{Supplied transport data}
    \label{subsec:transport}

    A comparison between fibres at neighbouring points $x$ and $x + dx$
    requires the specification of a transport map
    \begin{equation}
      T_{x \to x+dx} : F_x \;\longrightarrow\; F_{x+dx}.
    \end{equation}
    Because the projection $\Pi$ is non-injective, there is no canonical
    identification between fibres at distinct points, and the projection does not supply one.
    The differential of the bundle projection maps tangent vectors of the total space to tangent
    vectors of the base: it fixes which directions are vertical, but it does not choose horizontal
    lifts, and hence it does not identify neighbouring fibres.
    The two continuum hypotheses therefore do not determine a transport.
    On the trivial bundle $M \times \mathbb{C}$ over a coordinate patch, for instance, the covariant
    derivatives $d$ and $d + i\,x\,dy$ have the same projection and the same fibres, while the first is
    flat and the second has curvature $i\,dx \wedge dy$.

    The development below therefore rests on a third input, supplied and not derived:
    \begin{enumerate}[(T1)]
      \item a reduction of the fibre bundle to a principal bundle with a compact Lie structure group
      $G$, of Lie algebra $\mathfrak{g}$;
      \item a unitary representation of $G$ on the fibre, states being local sections $\psi$ of the
      associated vector bundle;
      \item a connection on that principal bundle, acting on $\psi$ through the representation.
    \end{enumerate}
    The structure group $G$ is distinct from the group $\mathcal{G}$ of
    Definition~\ref{def:gauge-group}, whose elements are global automorphisms of
    $\Omega_n^{(c)}$: nothing above reduces $\mathcal{G}$ to a finite-dimensional Lie group.
    The group read off the supplied carrier of Proposition~\ref{thm:su2-rigidity} is one possible
    choice of $G$, not one that this paper derives.
    A connection is the datum that assigns a holonomy to each path in the base and so defines
    parallel transport~\cite{BaezSchreiber2007}; the same bundle and projection admit many
    connections, as the example above shows.
    The data (T1)--(T3) might later be constructed from additional admissibility structure; no such
    construction is given here, and it is an open direction (Section~\ref{sec:open}).

    \emph{Conventions.}
    In a local trivialisation, write $D_\mu = \partial_\mu + A_\mu$, with $A_\mu$ anti-Hermitian
    and $\mathfrak{g}$-valued in the representation (T2).
    Parallel transport along a path $x(t)$ is defined by $\dot{x}^\mu D_\mu \psi = 0$, so that, to first
    order,
    \begin{equation}
      \label{eq:transport-expansion}
      T_{x \to x+dx}
      \;=\; \mathrm{Id} - A_\mu(x)\,dx^\mu + o(dx).
    \end{equation}
    Under a local gauge transformation $\psi \mapsto g\psi$ with $g(x) \in G$, the connection
    transforms as $A_\mu \mapsto g A_\mu g^{-1} - (\partial_\mu g)\,g^{-1}$; then
    $D_\mu\psi \mapsto g\,D_\mu\psi$ and $T_{x \to x+dx} \mapsto g(x+dx)\,T_{x \to x+dx}\,g(x)^{-1}$
    to first order.
    The connection is thus the infinitesimal generator of the supplied transport; it is not forced by
    the requirement of comparing admissible fibres.

  \subsection{Gauge connection from fibre transport}

    \begin{definition}[Gauge connection]
      \label{def:connection}
      On the supplied data (T1)--(T3), a \emph{gauge connection} $A_\mu$ is the $\mathfrak{g}$-valued
      one-form appearing in~\eqref{eq:transport-expansion}, defining a covariant derivative
      \begin{equation}
        \label{eq:covariant-derivative}
        D_\mu \;:=\; \partial_\mu + A_\mu
      \end{equation}
      such that $D_\mu\psi$ transforms homogeneously under local gauge
      transformations.
    \end{definition}

    \begin{remark}[Admissibility reading of $A_\mu$]
      \label{prop:connection-admissibility}
      In the projection framework, $A_\mu(x)$ is read as encoding the local variation of the
      admissible fibre structure: the infinitesimal correction that the supplied transport applies to
      an admissible configuration when it is carried from the fibre at $x$ to the fibre at $x + dx$.
      This is an interpretation, not a result.
      It requires that the supplied transport map admissible configurations to admissible
      configurations, which is a further condition on (T1)--(T3); it is not derived from the
      admissibility constraints (A1--A3 of~\cite{Beau2026m}).
      Those constraints define, at each point $x$, a fibre $F_x$ of structurally compatible
      configurations, but they neither select a connection nor single out a linear correction
      $A_\mu(x)\,dx^\mu$: the differential of the projection relates base directions, not fibres
      (§\ref{subsec:transport}).
    \end{remark}

    \begin{remark}[Relation to the holonomy argument]
        The white paper~\cite{Cosmochrony} shows that admissible injective projections carry
        trivial holonomy, so that non-trivial holonomy, the minimal structural condition for a
        gauge symmetry to arise as a descriptive invariance, requires non-injectivity; on a
        minimal model it illustrates how admissibility constraints produce a non-contractible
        loop in the space of admissible configurations.
        The reading of Remark~\ref{prop:connection-admissibility} is the local infinitesimal
        counterpart of this setting, on supplied transport data.
        Neither statement establishes non-trivial holonomy for a given non-injective
        projection.
    \end{remark}

  \section{Curvature and the Yang\textendash Mills Action}
  \label{sec:ym}

  \subsection{Curvature as path dependence of transport}

    Given the transport map of §\ref{subsec:transport}, consider two infinitesimal
    displacements $dx^\mu$ and $dx^\nu$ from a point $x$.
    The failure of successive transports along these directions to commute defines
    the curvature of the connection.

    \begin{definition}[Field strength]
      \label{def:field-strength}
      The \emph{field strength} associated with the gauge connection $A_\mu$ is the
      curvature of the covariant derivative:
      \begin{equation}
        \label{eq:field-strength}
        F_{\mu\nu} \;:=\; [D_\mu, D_\nu]
        \;=\; \partial_\mu A_\nu - \partial_\nu A_\mu + [A_\mu, A_\nu].
      \end{equation}
    \end{definition}

    $F_{\mu\nu}$ measures the path dependence of the supplied transport around infinitesimal loops:
    transport around the small parallelogram spanned by $dx^\mu$ and $dx^\nu$ differs from the identity
    by a term proportional to $F_{\mu\nu}\,dx^\mu dx^\nu$.
    Curvature is not a failure of admissibility.
    A curved connection transports configurations consistently along each path; what it does not
    provide is an identification of distant fibres independent of the path.
    The paper constructs no admissibility constraint operators, so it does not identify $F_{\mu\nu}$
    with a non-commutativity of admissibility constraints; that identification would require those
    operators and their relation to $D_\mu$, and it is open.

    Vanishing curvature has a local meaning only.
    A flat connection ($F_{\mu\nu} = 0$) gives a path-independent identification of fibres over a
    simply connected region; globally, it can still have a non-trivial holonomy around a
    non-contractible loop~\cite{TongGauge}.
    On the circle, for instance, the connection $A = i\alpha\,d\theta$ with a constant real $\alpha$ is
    flat; with the convention~\eqref{eq:transport-expansion}, its holonomy is $e^{-2\pi i\alpha}$, which
    is non-trivial for every non-integer $\alpha$.
    A flat connection therefore does not by itself give a globally path-independent fibre structure.

  \subsection{The Yang\textendash Mills action as a structural motivation}

    \paragraph{Structural motivation for Yang\textendash Mills-type dynamics.}
      The admissibility principle (A3 of~\cite{Beau2026m}) favours configurations
      that preserve the full multiplicity of open directions without premature
      selection.
      Translated to the gauge connection, this principle is read as favouring configurations
      that minimise the path dependence of fibre transport, i.e.\ minimise the
      curvature $\|F_{\mu\nu}\|^2$; this reading is a motivation, not a consequence of A3.

      On the supplied data (T1)--(T3), $F_{\mu\nu}$ is the curvature of the supplied connection.
      The paper does not derive $F_{\mu\nu}$ from the non-injectivity of $\Pi$, since the continuum
      hypotheses do not determine the connection (§\ref{subsec:transport}).
      What non-injectivity contributes is the setting: fibres without a canonical identification
      across points, in which a connection is needed; it does not supply the connection.

      The subsequent step is a standard result of gauge theory, not specific to the present
      framework, and it holds under an explicit truncation.
      On a fixed four-dimensional Lorentzian background, consider local Lagrangians that are
      gauge-invariant, Lorentz-invariant and parity-even, built from $F_{\mu\nu}$ and its covariant
      derivatives, of total mass dimension at most four, equivalently with at most two derivatives of
      $A_\mu$ in total.
      Apart from a constant, the only such term is an invariant quadratic form on $\mathfrak{g}$
      evaluated on $F_{\mu\nu}$ and $F^{\mu\nu}$.
      The truncation excludes higher powers such as $(F_{\mu\nu}F^{\mu\nu})^2$, whose field equations are
      of second order but which carry four derivatives in total.
      For a compact $\mathfrak{g}$ that is a sum of simple and abelian factors, the quadratic invariant
      is fixed up to one positive coefficient per simple factor and a positive-definite symmetric
      matrix on the abelian factors; that matrix, which allows kinetic mixing, becomes the identity in a
      suitable basis of the abelian factors.
      With these qualifications, the action is $\int\mathrm{Tr}(F_{\mu\nu}F^{\mu\nu})$ with one coupling
      per simple factor and, in that basis, one per abelian factor.

      \begin{remark}[Yang\textendash Mills action]
        \label{rem:ym-action}
        Under the continuum hypotheses of §\ref{sec:connection}, the minimisation
        principle above is read as favouring the action
        \begin{equation}
          \label{eq:ym-action}
          S_{\mathrm{YM}}[A]
          \;=\; \int \mathrm{Tr}\bigl(F_{\mu\nu} F^{\mu\nu}\bigr)\,\sqrt{-g}\,d^4x,
        \end{equation}
        whose variation yields the Yang\textendash Mills equations $D_\mu F^{\mu\nu} = 0$~\cite{YangMills1954}.
        Here $\mathrm{Tr}$ is the trace in the representation (T2), negative-definite on the
        anti-Hermitian algebra $\mathfrak{g}$; a positive coupling per factor is understood
        in~\eqref{eq:ym-action}, the metric $g$ is a fixed background, and the variation is taken over
        connections with compactly supported variations, so that boundary terms vanish.
        This is a structural motivation, not a derivation from Axioms~A1--A4 alone.
        The uniqueness of~\eqref{eq:ym-action} holds under the truncation stated above (total mass
        dimension at most four, gauge-invariant, Lorentz-invariant, parity-even; one coupling per
        simple factor and, in a basis that absorbs kinetic mixing, per abelian factor); it does not
        require the Cosmochrony framework.
        The framework does not supply $F_{\mu\nu}$ itself either: it is the curvature of the supplied
        connection of §\ref{subsec:transport}.
        The Born\textendash Infeld capacity paper~\cite{Beau2026c} does not supply this step:
        it leaves the relational construction of an antisymmetric gauge curvature as an open
        problem.
        A derivation of~\eqref{eq:ym-action} from Axioms~A1--A4 is open
        (Section~\ref{sec:open}).
      \end{remark}

  \section{Open Directions}
  \label{sec:open}

  \paragraph{Derivation of the Yang\textendash Mills action.}
    Remark~\ref{rem:ym-action} gives a structural motivation for the
    Yang\textendash Mills action.
    Beyond the standard power-counting argument, neither its uniqueness among
    admissibility-compatible actions nor its derivation from Axioms~A1--A4 is established.
    The relational construction of an antisymmetric curvature two-form is itself
    open~\cite{Beau2026c}.

  \paragraph{Transport data.}
    The continuum hypotheses of Section~\ref{sec:connection} do not determine a connection.
    Constructing a structure group, its representation and a connection, the data (T1)--(T3) of
    §\ref{subsec:transport}, from admissibility structure is open, and so is the condition that the
    resulting transport preserve admissible configurations.

  \paragraph{Physical charges.}
    Gauge charges in the sense of Definition~\ref{def:gauge-charge} are functions of the observable
    state.
    Identifying them with conserved physical charges requires a dynamics, a charge generator and a
    matter representation, none of which is given here.

  \paragraph{Carrier selection.}
    The $\mathrm{SU}(2)$ sector of Proposition~\ref{thm:su2-rigidity} rests on a supplied carrier
    and on the hypothesis $\mathrm{Im}\,\Pi_{\mathrm{adm}} = \mathrm{Ntrl}$.
    Selecting the carrier~\cite{Beau2026a27} and deriving that hypothesis~\cite{Beau2026a31} are
    open.

  \paragraph{$\mathrm{SU}(3)$ from admissibility.}
    The identification of $\mathrm{SU}(3)$ as a gauge group requires a carrier, supplied or
    derived, on which an $\mathrm{SU}(3)$-type group acts within the admissible fibre, together
    with a selection argument; spectral degeneracies alone do not supply it.

  \paragraph{Spatial continuum.}
    The second continuum hypothesis of Section~\ref{sec:connection} is the open spatial
    continuum question of~\cite{Beau2026q5a}.

  \paragraph{Coupling constants.}
    The present paper does not determine the numerical values of coupling constants
    ($g_s$, $g_w$, $\alpha$), and no result of the programme determines them.
    In particular, the quantity $\beta^*$ of the spectral programme has no native carrier on the
    Heisenberg substrate~\cite{Beau2026a31}, so no transfer from spectral data to a coupling is
    available.

  \section{Conclusion}
  \label{sec:conclusion}

    The paper defines gauge charges as functions constant on the fibres of a non-injective
    projection, and the gauge group as the group of fibre-preserving transformations.
    A gauge charge factors through the projection, so it is a function of the observable state;
    its identification with a conserved physical charge is open.
    $\mathrm{U}(1)$ is the residual phase symmetry of span-based observables, and $\mathrm{SU}(2)$
    is read off a supplied carrier under a neutral-sector hypothesis; neither result selects a gauge
    group, and $\mathrm{SU}(3)$ is a hypothesis.
    The two continuum hypotheses do not determine a connection, so a structure group, its
    representation and a connection are supplied.
    On these data the connection generates parallel transport, its curvature measures path
    dependence around infinitesimal loops, and flatness gives a path-independent identification of
    fibres only locally.
    The Yang\textendash Mills action is singled out by a standard truncation argument on these data
    and is motivated, not derived, by the admissibility principle.

  \subsection*{Interpretive outlook}

    What follows is a structural reading, offered as interpretation and not as a result.

    A non-injective projection leaves a fibre that the observables do not resolve, and the
    symmetries of that fibre are the natural place to read gauge redundancy.
    The charge result points in the expected direction for that reading: what is invariant under
    the redundancy is a function of the observable output, so gauge-invariant quantities are exactly
    those the observables already see.
    Gauge physics needs more than redundancy.
    It needs a particular group acting on the fibre, a rule for comparing fibres at different points,
    and a dynamics for that rule.
    The paper locates where each of these would enter and shows that none follows from
    non-injectivity together with a continuum description: the group is read off a supplied carrier,
    the comparison is a supplied connection, and the dynamics is a standard truncation on it.
    A third party should situate the paper as a dictionary between non-injective projection and
    gauge redundancy, and as a list of the inputs that a derivation of gauge fields would still have
    to supply, not as such a derivation.

  \section*{Acknowledgements}

  The author acknowledges the use of large language models as a supportive tool for
  refining language, structure, and internal consistency during the development of
  this manuscript.
  All conceptual contributions, theoretical choices, and interpretations remain the
  sole responsibility of the author.

  \clearpage
  \bibliographystyle{plain}
      \bibliography{cosmochrony-bibliography,external-refs}

\end{document}
